\documentclass[11pt]{article}
\usepackage{fontspec}
\setmainfont{Latin Modern Roman}
\usepackage{amsmath,amssymb,mathtools,unicode-math}
\setmathfont{Latin Modern Math}
\usepackage{xurl,hyperref}
\hypersetup{hidelinks,pdftitle={Problem statement}}
\newcommand{\sref}[2]{\hyperlink{source:#1:#2}{[S#2]}}
\newcommand{\cataloguescope}[1]{\paragraph{Mathematical question.}#1}
\begin{document}
\section{Research Problems in Function Theory — Problem 6.35}\label{statement}

\subsection{Definitions and mathematical statement}

\paragraph{Shared notation.}$\mathbb N=\{1,2,\ldots\}$, and $\mathbb Z,\mathbb Q,\mathbb R,\mathbb C$ denote the integers, rationals, reals and complex numbers. Any other conventions are specified locally.


Write $\mathbb D=\{z\in\mathbb C:|z|<1\}$ and $\mathbb T=\partial\mathbb D$. Let $S$ consist of holomorphic injective functions $f(z)=z+\sum_{n\ge2}a_nz^n$ on $\mathbb D$; set $a_0=0$, $a_1=1$. A function in $S$ is starlike if its image is starlike about $0$, and convex if its image is a convex domain.For $f\in S$, let $O_f=\mathbb D\setminus f(\mathbb D)$. Area means two-dimensional Lebesgue measure. A finite omitted configuration is a list $w_1,\ldots,w_m\in\mathbb D$, with $m\ge1$.
\paragraph{Statement recorded in the supplied source.}
Characterize the subsets $O\subset\mathbb D$ that equal $O_f$ for some $f\in S$; determine $\sup_{f\in S}\operatorname{area}(O_f)$ and the maximizing maps; and give necessary and sufficient conditions on any finite configuration for the existence of $f\in S$ with all $w_j\in O_f$. Include the two-point and three-point cases explicitly in that classification. \sref{statement}{2}
\subsection{Short English statement}
Classify the omitted parts of the unit disc for normalized univalent maps, maximize their area, and decide which finite point configurations can be omitted together.
\subsection{Sources}
\begin{description}
\item[\hypertarget{source:statement:1}{S1}] ULAM.AI and the UnsolvedMath Contributors, UnsolvedMath: A Curated Collection of Open Mathematics Problems, catalogue entry. CC BY 4.0. \url{https://huggingface.co/datasets/ulamai/UnsolvedMath}.
\item[\hypertarget{source:statement:2}{S2}] W. K. Hayman and E. F. Lingham, Research Problems in Function Theory (2018), arXiv:1809.07200, Chapter 6, Problem 6.35, complete update and preceding notation. \url{https://arxiv.org/abs/1809.07200}.
\end{description}
\label{end:statement}
\end{document}
